\documentclass[11pt]{article}

\usepackage{amssymb,amsmath,amsfonts,amsthm}
\usepackage{lmodern}
\usepackage{hyperref}
\usepackage{booktabs}

\setlength{\textwidth}{15.5cm} \setlength{\headheight}{0.5cm} \setlength{\textheight}{21.5cm}
\setlength{\oddsidemargin}{0.25cm} \setlength{\evensidemargin}{0.25cm} \setlength{\topskip}{0.5cm}
\setlength{\footskip}{1.5cm} \setlength{\headsep}{0cm} \setlength{\topmargin}{0.5cm}

\newcommand{\Z}{\mathbb{Z}}
\newcommand{\R}{\mathbb{R}}
\newcommand{\pbar}{\bar{p}}
\newcommand{\tw}{t_{\mathrm{w}}}

\newtheorem{definition}{Definition}
\newtheorem{proposition}{Proposition}
\newtheorem{remark}{Remark}

\begin{document}

\section*{QP Formulation for Marked Edge Selection}

\subsection*{1.\ Setup}

Let $G = (V, E)$ be a connected graph with $|V| = n$, populations $\{p_v\}$,
$d$ districts under the current partition $\xi$, and ideal district population
$\pbar = \frac{1}{d}\sum_v p_v$. Let $T'$ be a rooted spanning tree with
$n - 1$ edges.

For each edge $e_i$ (connecting child node $v_i$ to its parent), define:
\begin{itemize}
  \item \textbf{Subtree population:}
    $\displaystyle \sigma_i = \sum_{u \in \mathrm{subtree}(v_i)} p_u$.
  \item \textbf{Ancestor indicator:}
    $A_{ij} = 1$ if $e_j$ is a descendant of $e_i$ in $T'$, and $0$ otherwise.
\end{itemize}

\subsection*{2.\ Divergence-Driven Flow}

The partition $\xi$ induces a population imbalance across districts. We encode
this imbalance as a source/sink vector and solve for the resulting flow on the
tree.

\subsubsection*{Source vector}

For each node $v \in V$ with $\xi(v) = i$, define the \textbf{district
imbalance source}:
\[
  b_v = \frac{p(\xi_i) - \pbar}{|\xi_i|},
\]
where $p(\xi_i) = \sum_{u \in \xi_i} p_u$ is the total population of
district~$\xi_i$ and $|\xi_i|$ is its number of nodes. Nodes in overpopulated
districts ($p(\xi_i) > \pbar$) act as sources; nodes in underpopulated
districts act as sinks. By construction,
\[
  \sum_{v \in V} b_v = \sum_{i=1}^d \bigl(p(\xi_i) - \pbar\bigr) = 0,
\]
satisfying the solvability condition for the Laplacian.

\subsubsection*{Population potential}

Solve the Poisson equation on the graph:
\[
  \Delta_G\, \phi = b.
\]
The solution $\phi \in \R^V$ is the \textbf{population potential}. For
planar redistricting graphs, the sparse Cholesky factorization of $\Delta_G$
costs $O(n^{3/2})$; each back-solve is $O(n)$.

\subsubsection*{Tree flow}

On a tree, the flow on each edge is uniquely determined by $b$ (no cycles).
For each tree edge $e_i$ (child $v_i$ to parent), the
\textbf{rebalancing flow} is:
\[
  f_i = \sum_{u \in \mathrm{subtree}(v_i)} b_u
  = \sum_{u \in \mathrm{subtree}(v_i)}
    \frac{p(\xi_{\xi(u)}) - \pbar}{|\xi_{\xi(u)}|}.
\]
This measures the net population pressure crossing edge $e_i$: large $|f_i|$
means $e_i$ channels significant rebalancing flow between surplus and deficit
districts.

\begin{remark}[Flow vs.\ subtree population]
The subtree population $\sigma_i$ is a static property of the tree --- it does
not depend on the current partition $\xi$. The flow $f_i$ is
\textbf{partition-aware}: it reflects the inter-district imbalance. Two edges
with similar $\sigma_i$ values may have very different $f_i$ values if one
sits between two balanced districts ($f_i \approx 0$) while the other channels
pressure between a surplus and deficit district ($|f_i|$ large).
\end{remark}


\subsection*{3.\ The QP Objective}

Selecting $S \subset E(T')$ with $|S| = d - 1$ edges to remove creates
$d$ components. The \textbf{total population imbalance} is
\[
  F(S) = \sum_{\text{components } c} \bigl(\mathrm{pop}(c) - \pbar\bigr)^2.
\]
We combine the structural balance criterion ($\sigma_i$) with the
flow-based pressure criterion ($f_i$) into a single objective.

\subsubsection*{Linear coefficients (combined)}

For each tree edge $e_i$, the linear term balances two goals:
\[
  c_i = \underbrace{(\sigma_i - \pbar)^2 + (n - \sigma_i - \pbar)^2
    - (n - \pbar)^2}_{\text{population balance (cut quality)}}
  \;-\; \lambda\, f_i^2,
\]
where:
\begin{itemize}
  \item The first term penalizes edges whose removal creates unbalanced
    components (small $|\sigma_i - \pbar|$ is preferred).
  \item The second term $-\lambda\, f_i^2$ \textbf{rewards} edges under high
    rebalancing pressure (large $|f_i|$ means the edge channels flow from
    surplus to deficit).
  \item The parameter $\lambda > 0$ controls the trade-off. In practice,
    $\lambda = 1$ or $\lambda$ chosen so the two terms have comparable
    magnitude.
\end{itemize}

\subsubsection*{Pairwise interaction coefficients}

The interaction $Q_{ij}$ encodes the ancestor--descendant dependency. When two
selected edges are nested, the component between them has population
$\sigma_i - \sigma_j$ (not $\sigma_i$ or $\sigma_j$ independently).

Using inclusion--exclusion:
$Q_{ij} = F(\{e_i, e_j\}) - F(\{e_i\}) - F(\{e_j\}) + F(\emptyset)$, giving:

\[
  Q_{ij} =
  \begin{cases}
    0 & \text{if } i = j, \\[4pt]
    2\,\sigma_j\,(n - \sigma_i)
      & \text{if } A_{ij} = 1
        \text{ ($e_j$ descendant of $e_i$, $\sigma_j \leq \sigma_i$)}, \\[4pt]
    2\,\sigma_i\,(n - \sigma_j)
      & \text{if } A_{ji} = 1
        \text{ ($e_i$ descendant of $e_j$, $\sigma_i \leq \sigma_j$)}, \\[4pt]
    2\,\sigma_i\,\sigma_j
      & \text{otherwise (non-ancestor pair).}
  \end{cases}
\]

\begin{remark}[Dependency encoding]
For ancestor--descendant pairs, $Q_{ij} = 2\,\sigma_j(n - \sigma_i) <
2\,\sigma_i\,\sigma_j$ whenever $\sigma_i > n/2$. The QP
\textbf{discourages} selecting nested edges with large overlapping subtrees.
For non-ancestor pairs, $Q_{ij} = 2\,\sigma_i\,\sigma_j$ penalizes selecting
two edges that both cut off large subtrees (their ``remainder'' component
$n - \sigma_i - \sigma_j$ would be too small).
\end{remark}


\subsection*{4.\ The QP}

\[
\boxed{
  \min_{x \in [0,1]^{n-1}} \; x^\top Q\, x + c^\top x
  \quad \text{subject to} \quad
  \sum_{i=1}^{n-1} x_i = d - 1
}
\]

After solving the continuous relaxation, round by selecting the $d - 1$ edges
with the largest $x_i$ values. Check population balance of the resulting
partition $\xi' = T' \setminus M'$.

\subsection*{5.\ Exactness}

\begin{proposition}
The pairwise QP (the $Q_{ij}$ terms) is \textbf{exact} for $d \leq 3$
(selecting $d - 1 \leq 2$ edges). For $d \geq 4$, there exist three-body
and higher-order corrections:
\[
  F(S) = F_{\mathrm{QP}}(x) + \sum_{i < j < k} R_{ijk}\, x_i\, x_j\, x_k
  + \cdots
\]
The three-body residual $R_{ijk}$ is nonzero only when $e_i, e_j, e_k$ form
an ancestor chain (all three lie on a single root-to-leaf path). For balanced
partitions, such deep nesting is rare, and the pairwise approximation is
accurate.
\end{proposition}


\subsection*{6.\ Summary of the Combined Approach}

The QP integrates two sources of information:

\medskip
\begin{center}
\begin{tabular}{lll}
\toprule
\textbf{Component} & \textbf{Source} & \textbf{Role} \\
\midrule
$\sigma_i$ & Tree structure & Cut quality (is $\sigma_i \approx \pbar$?) \\
$f_i$ & Partition $\xi$ + Laplacian & Rebalancing pressure (where is flow?) \\
$Q_{ij}$ & Ancestor--descendant relation & Dependency between nested edges \\
\bottomrule
\end{tabular}
\end{center}

\medskip\noindent
\textbf{Computational cost per MCMC step:}
\begin{itemize}
  \item Subtree populations $\sigma_i$: $O(n)$ via DFS.
  \item Source vector $b$ and tree flow $f_i$: $O(n)$ via DFS (given the
    current partition~$\xi$).
  \item Population potential $\phi$: $O(n)$ back-solve (reusing the Cholesky
    factor of $\Delta_G$, pre-computed in $O(n^{3/2})$ for planar graphs).
  \item QP solve: $O(n)$ for the continuous relaxation with the equality
    constraint (closed-form via Lagrange multiplier on $\sum x_i = d - 1$),
    followed by rounding.
  \item Total per step: $O(n)$, dominated by the tree traversals.
\end{itemize}

\end{document}
